% CVPR 2023 Paper Template
% based on the CVPR template provided by Ming-Ming Cheng (https://github.com/MCG-NKU/CVPR_Template)
% modified and extended by Stefan Roth (stefan.roth@NOSPAMtu-darmstadt.de)

\documentclass[10pt,twocolumn,letterpaper]{article}

%%%%%%%%% PAPER TYPE  - PLEASE UPDATE FOR FINAL VERSION
\usepackage{cvpr}      % To produce the REVIEW version
%\usepackage{cvpr}              % To produce the CAMERA-READY version
%\usepackage[pagenumbers]{cvpr} % To force page numbers, e.g. for an arXiv version

% Include other packages here, before hyperref.
\usepackage{graphicx}
\usepackage{amsmath}
\usepackage{amssymb}
\usepackage{booktabs}
\usepackage{subcaption}


% It is strongly recommended to use hyperref, especially for the review version.
% hyperref with option pagebackref eases the reviewers' job.
% Please disable hyperref *only* if you encounter grave issues, e.g. with the
% file validation for the camera-ready version.
%
% If you comment hyperref and then uncomment it, you should delete
% ReviewTempalte.aux before re-running LaTeX.
% (Or just hit 'q' on the first LaTeX run, let it finish, and you
%  should be clear).
\usepackage[pagebackref,breaklinks,colorlinks]{hyperref}


% Support for easy cross-referencing
\usepackage[capitalize]{cleveref}
\crefname{section}{Sec.}{Secs.}
\Crefname{section}{Section}{Sections}
\Crefname{table}{Table}{Tables}
\crefname{table}{Tab.}{Tabs.}


%%%%%%%%% PAPER ID  - PLEASE UPDATE
\def\cvprPaperID{N/A} % *** Enter the CVPR Paper ID here
\def\confName{CVPR}
\def\confYear{2023}


\begin{document}

%%%%%%%%% TITLE - PLEASE UPDATE
\title{Plant Disease Recognition via Transfer Learning and Explainability on PlantVillage}

\author{
Yang Zhang\\
Concordia University\\
{\tt\small 40332950}
\and
Yongqi Hao\\
Concordia University\\
{\tt\small 40303141}
\and
Siming Yi\\
Concordia University\\
{\tt\small 40333494}
\and
Wenbin Zhang\\
Concordia University\\
{\tt\small 40332403}
}

\maketitle

\begin{abstract}
Accurate plant disease recognition from leaf images is important for precision agriculture and early crop intervention. In this project, we study image-based plant disease classification on the PlantVillage dataset, which contains 54,305 images across 38 classes and 14 plant species. We compare transfer learning strategies and model families, including ResNet-50 trained from scratch, EfficientNet-B3 with linear probing and full fine-tuning, and ViT-Small with full fine-tuning. We also analyze the effect of input modality using color, grayscale, and background-segmented images. Our results show that transfer learning is essential for this task: full fine-tuning of pretrained models greatly outperforms both linear probing and from-scratch training. EfficientNet-B3 achieves the best validation Macro-F1 of 99.28\% on color images, while ViT-Small achieves 99.34\% test accuracy and 99.09\% test Macro-F1 among the fully evaluated runs. An input-modality ablation further shows that color is a critical diagnostic cue, whereas background removal provides only marginal benefit. Finally, GradCAM and GradCAM++ visualizations indicate that the model generally attends to lesion-bearing regions, with remaining errors concentrated in visually similar disease classes.
\end{abstract}

\section{Introduction}

Plant diseases significantly reduce agricultural productivity and can spread rapidly if not detected early. Manual diagnosis by domain experts is often costly and difficult to scale, especially when monitoring large crop areas. Computer vision offers a practical alternative by enabling automatic recognition of disease symptoms directly from leaf images.

Deep learning has substantially improved image classification performance, but plant disease recognition remains challenging in fine-grained multi-class settings. Different diseases may share similar textures, discoloration patterns, or lesion shapes, particularly within the same crop family. In addition, it is not always clear whether classification performance depends mainly on leaf symptoms, global color statistics, or background artifacts.

In this project, we investigate three questions: (1) how much transfer learning matters for PlantVillage classification, (2) whether removing the background improves performance, and (3) whether model predictions can be interpreted through saliency-based explainability. To answer these questions, we build a unified training pipeline, compare multiple architectures and training strategies, and analyze learned attention using GradCAM and GradCAM++. This framing is useful not only for achieving strong benchmark performance, but also for understanding what kinds of visual evidence modern models rely on in a relatively clean agricultural dataset. As a result, the project serves both as a classification study and as an analysis of representation quality, modality dependence, and trustworthiness of model behavior.

Our code is available at:
\url{https://github.com/Winter-ssolstice/COMP6341_Project}

This project was conducted solely for COMP6341 and is not combined with any other course project.

\section{Related Work}

PlantVillage has become a standard benchmark for plant disease classification since the dataset was introduced by Hughes and Salath\'e~\cite{hughes2015open}. Early work on this benchmark and related tasks mainly relied on convolutional neural networks and transfer learning from ImageNet-pretrained models. More recently, transformer-based architectures have shown strong performance on image recognition by capturing longer-range dependencies and richer global representations~\cite{dosovitskiy2021vit,tan2019efficientnet,he2016resnet}.

Explainability is also important in vision-based diagnosis. Gradient-based localization methods such as GradCAM~\cite{selvaraju2017grad} and GradCAM++~\cite{chattopadhay2018gradcampp} provide a way to visualize which image regions contribute most strongly to a prediction, helping assess whether a model focuses on disease-relevant evidence. In practical agricultural applications, this is especially useful because high classification accuracy alone may not be sufficient for user trust; growers and researchers often want visual evidence that the model responds to lesion regions rather than spurious background cues or dataset-specific artifacts.



Beyond PlantVillage, CNN-based approaches have been widely adopted for plant disease recognition because they can capture local textures and lesion morphology directly from leaf images. In practical deployments, transfer learning is often preferred since pretrained visual representations provide a strong initialization and improve data efficiency. In addition, explainability has become increasingly important in visual diagnosis tasks: high accuracy alone is not always sufficient without evidence that a model attends to semantically meaningful regions. GradCAM-style methods provide a simple and widely used mechanism for this purpose by highlighting image areas that most influence the predicted class.

\section{Approach}

\subsection{Dataset and Preprocessing}

We use the PlantVillage dataset, which contains 54,305 leaf images from 38 disease/health classes across 14 plant species. Three input versions are considered: \textit{color}, \textit{grayscale}, and \textit{background-segmented}. All images are resized to $224 \times 224$. We use a fixed 80/10/10 split on the original dataset, resulting in 43,444 training, 5,430 validation, and 5,431 test images for the color version. The grayscale version reuses the same image identities. The background-segmented version follows the same source split when segmented counterparts are available; in the aligned evaluation used later, 5,318 of the 5,431 test images have segmented matches.

To improve generalization, we apply data augmentation during training: random horizontal flip, random vertical flip, random resized crop, color jitter, and MixUp. All models are optimized with cross-entropy loss using AdamW and cosine learning-rate scheduling. Model selection is based on validation performance, and the best checkpoint for each experiment is used for final evaluation. This unified protocol ensures that differences across runs mainly reflect model design and transfer-learning strategy rather than inconsistent split construction or evaluation settings.

All experiments were implemented in PyTorch. Model selection was based on validation performance, and the checkpoint with the best validation result was used for final reporting. Batch size and total training epochs were chosen according to the memory and convergence characteristics of each architecture, while keeping the data split and evaluation metrics fixed across settings.

\subsection{Model Comparison}

We compare four model-strategy settings on the color dataset:
\begin{itemize}
    \item \textbf{ResNet-50 from scratch}: no pretrained weights;
    \item \textbf{EfficientNet-B3 linear probing}: pretrained backbone frozen, classifier head trained;
    \item \textbf{EfficientNet-B3 full fine-tuning}: pretrained model with all layers updated;
    \item \textbf{ViT-Small full fine-tuning}: pretrained transformer with all layers updated.
\end{itemize}

This setup isolates the effects of transfer learning and architecture choice. We report validation and test metrics using Top-1 accuracy and Macro-F1.

\subsection{Input-Modality Ablation}

After model comparison, we analyze the effect of input modality. We evaluate the selected strong pretrained transformer pipeline on the three dataset versions: color, grayscale, and background-segmented. This allows us to measure the contribution of color information and test whether removing the background simplifies the task.

\subsection{Explainability}

To interpret model behavior, we apply GradCAM and GradCAM++ to representative correct and incorrect predictions. The goal is to verify whether the model focuses on lesion-bearing regions and to diagnose failure modes in confusing classes.

\section{Experiments}

\begin{figure*}[!t]
    \centering
    \begin{subfigure}{0.46\linewidth}
        \centering
        \includegraphics[width=\linewidth]{part3/vit_small_full_finetune/comparisons/correct_tomato___septoria_leaf_spot.png}
        \caption{Correct prediction: Tomato / Septoria leaf spot.}
        \label{fig:gradcam-correct}
    \end{subfigure}
    \hfill
    \begin{subfigure}{0.46\linewidth}
        \centering
        \includegraphics[width=\linewidth]{part3/vit_small_full_finetune/comparisons/incorrect_tomato___spider_mites_two_spotted_spider_mite.png}
        \caption{Incorrect prediction: Tomato / Spider mites two-spotted spider mite.}
        \label{fig:gradcam-wrong}
    \end{subfigure}
    \caption{\textbf{GradCAM and GradCAM++ visualizations for representative correct and incorrect predictions.}
    For the correct case, the model attends to lesion-bearing regions. For the incorrect case, the attention becomes more diffuse and less discriminative. The segmented variants show similar behavior, consistent with the small performance gap between color and background-segmented inputs.}
    \label{fig:gradcam-results}
\end{figure*}


\subsection{Baseline}

We first establish a baseline using a pretrained ResNet-18 model fully fine-tuned on the color dataset. This baseline achieves 97.53\% test accuracy and 96.34\% test Macro-F1, showing that standard convolutional transfer learning is already highly effective on PlantVillage.

\subsection{Model Comparison}

\Cref{tab:model-selection} summarizes validation-set model selection on the color dataset. The most important pattern is that transfer learning is indispensable: the from-scratch ResNet-50 baseline performs poorly, and linear probing is also far below full fine-tuning. EfficientNet-B3 with full fine-tuning achieves the best validation Macro-F1 (99.28\%), with ViT-Small close behind at 98.89\%. This indicates that generic pretrained features alone are not sufficient; the backbone must be adapted to the target domain.

\Cref{tab:test-results} reports held-out test metrics only for runs with exported test evaluations. The pretrained ResNet-18 baseline already performs strongly, and ViT-Small with full fine-tuning reaches 99.34\% test accuracy and 99.09\% test Macro-F1 among the advanced models with complete held-out evaluation. For the later ablation and explainability analysis, we use the ViT-Small pipeline because it produced complete test artifacts and was the model used in our GradCAM workflow.

\begin{table}[t]
\centering
\small
\begin{tabular}{lcc}
\toprule
Model / Strategy & Val. Acc. & Val. Macro-F1 \\
\midrule
ResNet-50 from scratch (val) & 48.10 & 35.64 \\
EffNet-B3 linear probing & 67.31 & 58.01 \\
EffNet-B3 full fine-tune (val) & 99.54 & 99.28 \\
ViT-Small full fine-tune & 99.30 & 98.89 \\
\bottomrule
\end{tabular}
\caption{Validation-set model selection on the color dataset.}
\label{tab:model-selection}
\end{table}

\begin{table}[t]
\centering
\small
\begin{tabular}{lcc}
\toprule
Model / Strategy & Test Acc. & Test Macro-F1 \\
\midrule
ResNet-18 baseline & 97.53 & 96.34 \\
EffNet-B3 linear probing & 67.41 & 58.50 \\
ViT-Small full fine-tune & 99.34 & 99.09 \\
\bottomrule
\end{tabular}
\caption{Held-out test results for runs with exported test metrics.}
\label{tab:test-results}
\end{table}



These results suggest that PlantVillage is not merely an easy benchmark when evaluated under mismatched training strategies. Although pretrained models achieve near-saturated performance, the large gap between full fine-tuning and linear probing indicates that disease recognition still requires substantial domain adaptation. The poor from-scratch performance of ResNet-50 further highlights the importance of strong initialization and data-efficient optimization for fine-grained multi-class recognition.

\subsection{Input-Modality Ablation}

\Cref{tab:ablation} shows the effect of input modality using the ViT-Small evaluation pipeline. Color images produce the best results. Converting the images to grayscale causes the largest drop, reducing test accuracy to 95.86\% and Macro-F1 to 94.46\%. This suggests that chromatic information is a major diagnostic cue for PlantVillage diseases. Background segmentation, by contrast, changes the result only slightly: on the aligned subset with 5,318 matched test images, accuracy is 99.25\% and Macro-F1 is 99.09\%. Because this background-segmented score is computed on a slightly smaller matched subset than the color and grayscale rows, it should be interpreted as a near-matched comparison rather than a perfectly identical test-set evaluation.

\begin{table}[t]
\centering
\small
\begin{tabular}{lcc}
\toprule
Input Version & Test Acc. & Test Macro-F1 \\
\midrule
Color & 99.34 & 99.09 \\
Grayscale & 95.86 & 94.46 \\
Background-segmented & 99.25 & 99.09 \\
\bottomrule
\end{tabular}
\caption{Input-modality ablation using the selected ViT-Small evaluation pipeline. The background-segmented row is computed on the 5,318 matched test images available after segmentation.}
\label{tab:ablation}
\end{table}



Even with the matched-subset caveat, the small gap between the color and background-segmented settings suggests that the benchmark is largely leaf-centric and that the classifier does not depend heavily on background context. By contrast, the substantial drop under grayscale conversion indicates that many disease categories are characterized not only by shape and texture but also by subtle color cues such as yellowing, browning, or rust-like pigmentation.


\subsection{Explainability and Failure Analysis}

As shown in \Cref{fig:gradcam-results}, GradCAM and GradCAM++ visualizations indicate that the model generally attends to lesion-bearing regions for correct predictions rather than irrelevant image areas. In the correctly classified tomato example, the highlighted regions broadly align with visible disease patterns, supporting the plausibility of the learned decision process.

For incorrect predictions, the attention becomes less discriminative and more diffuse, often spreading across larger leaf regions without isolating the most diagnostic symptoms. This behavior is consistent with our error analysis, where visually similar disease classes remain the most challenging. In grayscale, the failure modes become more severe, especially for tomato diseases such as \textit{Early blight}, \textit{Target Spot}, and \textit{Spider mites}, further confirming that color is an important discriminative cue. An additional observation from \Cref{fig:gradcam-results} is that the segmented inputs produce attention maps that are visually similar to those of the color inputs, which aligns with the quantitative ablation where background removal yields only marginal changes. Therefore, the primary remaining challenge is class ambiguity among diseases with similar visual symptoms rather than background clutter. From a practical perspective, this also suggests that future improvements are more likely to come from better fine-grained representation learning or harder example mining than from further preprocessing of the already-clean PlantVillage backgrounds. Another implication is that explainability can help distinguish between meaningful confidence and accidental shortcut behavior. When the heatmaps align with lesion structure, the prediction is easier to trust; when the response is broad or unstable, the model output should be treated more cautiously. This type of qualitative inspection is especially valuable in agricultural applications, where deployment decisions may depend not only on top-line accuracy but also on whether the model behaves consistently across disease presentations.

\section{Conclusion}

We presented a systematic study of plant disease recognition on PlantVillage using transfer learning, modality ablation, and explainability analysis. Our experiments show that full fine-tuning of pretrained models clearly outperforms both linear probing and from-scratch training. EfficientNet-B3 achieved the best validation Macro-F1 on the color dataset, while ViT-Small achieved 99.34\% test accuracy and 99.09\% test Macro-F1 among the fully evaluated runs. The modality ablation demonstrated that color is a critical signal, whereas background removal provides only marginal benefit. Finally, GradCAM-based analysis showed that the model usually focuses on lesion-relevant regions, with errors concentrated in visually similar disease categories. Taken together, these findings suggest that PlantVillage is best viewed not simply as an easy benchmark but as a useful testbed for understanding how transfer learning, visual modality, and interpretability interact in fine-grained agricultural recognition. Future work could explore class-balanced optimization, stronger augmentation for confusing classes, and evaluation on less curated real-world datasets with more realistic lighting, occlusion, and background variation. More broadly, our results highlight that even when overall accuracy is already very high, careful comparison of training strategy and input representation still reveals meaningful differences in model behavior. This makes the benchmark a useful educational setting for studying both performance and reliability in computer vision systems.

{\small
\bibliographystyle{ieee_fullname}
\bibliography{egbib}
}

\end{document}
